\documentclass[12pt]{article}
\usepackage{amsmath,amssymb,amsfonts}
\usepackage[margin=1in]{geometry}

\begin{document}

\begin{center}
{\large\bfseries 2nd Irish Mathematical Olympiad\\
29 April 1989}
\end{center}

\bigskip
\noindent\textbf{Paper 1}
\bigskip

\begin{enumerate}
\item A quadrilateral $ABCD$ is inscribed, as shown, in a square of area one unit. Prove that \[2 \le |AB|^2 + |BC|^2 + |CD|^2 + |DA|^2 \le 4.\]

\item A $3 \times 3$ magic square, with magic number $m$, is a $3 \times 3$ matrix such that the entries on each row, each column and each diagonal sum to $m$. Show that if the square has positive integer entries, then $m$ is divisible by 3, and each entry of the square is at most $2n - 1$, where $m = 3n$. [An example of a magic square with $m = 6$ is $\begin{pmatrix} 2 & 1 & 3 \\ 3 & 2 & 1 \\ 1 & 3 & 2 \end{pmatrix}$.]

\item A function $f$ is defined on the natural numbers $\mathbb{N}$ and satisfies the following rules: \begin{enumerate} \item[(a)] $f(1) = 1$; \item[(b)] $f(2n) = f(n)$ and $f(2n+1) = f(2n) + 1$ for all $n \in \mathbb{N}$. \end{enumerate} Calculate the maximum value $m$ of the set $\{f(n) : n \in \mathbb{N}, 1 \le n \le 1989\}$, and determine the number of natural numbers $n$, with $1 \le n \le 1989$, that satisfy the equation $f(n) = m$.

\item Note that $12^2 = 144$ end in two 4's and $38^2 = 1444$ end in three 4's. Determine the length of the longest string of equal nonzero digits in which the square of an integer can end.

\item Let $x = a_1a_2\ldots a_n$ be an $n$-digit number, where $a_1, a_2, \ldots, a_n$ ($a_1 \ne 0$) are the digits. The $n$ numbers \[x_1 = x = a_1a_2\ldots a_n, \quad x_2 = a_na_1\ldots a_{n-1}, \quad x_3 = a_{n-1}a_na_1\ldots a_{n-2},\] \[x_4 = a_{n-2}a_{n-1}a_na_1\ldots a_{n-3}, \quad \ldots, \quad x_n = a_2a_3\ldots a_na_1\] are said to be obtained from $x$ by the cyclic permutation of digits. Find, with proof, (i) the smallest natural number $n$ for which there exists an $n$-digit number $x$ such that the $n$ numbers obtained from $x$ by the cyclic permutation of digits are all divisible by 1989; and (ii) the smallest natural number $x$ with this property.

\end{enumerate}

\newpage

\begin{center}
{\large\bfseries 2nd Irish Mathematical Olympiad\\
29 April 1989}
\end{center}

\bigskip
\noindent\textbf{Paper 2}
\bigskip

\begin{enumerate}
\setcounter{enumi}{5}
\item Suppose $L$ is a fixed line, and $A$ a fixed point not on $L$. Let $k$ be a fixed nonzero real number. For $P$ a point on $L$, let $Q$ be a point on the line $AP$ with $|AP|\cdot|AQ| = k^2$. Determine the locus of $Q$ as $P$ varies along the line $L$.

\item Each of the $n$ members of a club is given a different item of information. They are allowed to share the information, but, for security reasons, only in the following way: A pair may communicate by telephone. During a telephone call only one member may speak. The member who speaks may tell the other member all the information s(he) knows. Determine the minimal number of phone calls that are required to convey all the information to each other.

\item Suppose $P$ is a point in the interior of a triangle $ABC$, that $x, y, z$ are the distances from $P$ to $A, B, C$, respectively, and that $p, q, r$ are the perpendicular distances from $P$ to the sides $BC, CA, AB$, respectively. Prove that \[xyz \ge 8pqr,\] with equality implying that the triangle $ABC$ is equilateral.

\item Let $a$ be a positive real number, and let \[b = \sqrt[3]{a + \sqrt{a^2+1}} + \sqrt[3]{a - \sqrt{a^2+1}}.\] Prove that $b$ is a positive integer if, and only if, $a$ is a positive integer of the form $\frac{1}{2}n(n^2+3)$, for some positive integer $n$.

\item \begin{enumerate} \item[(i)] Prove that if $n$ is a positive integer, then \[\binom{2n}{n} = \frac{(2n)!}{(n!)^2}\] is a positive integer that is divisible by all prime numbers $p$ with $n < p \le 2n$, and that $\binom{2n}{n} < 2^{2n}$. \item[(ii)] For $x$ a positive real number, let $\pi(x)$ denote the number of prime numbers $p \le x$. [Thus, $\pi(10) = 4$ since there are 4 primes, viz., 2, 3, 5 and 7, not exceeding 10.] Prove that if $n \ge 3$ is an integer, then \begin{enumerate} \item[(a)] $\pi(2n) < \pi(n) + \dfrac{2n}{\log_2(n)}$; \quad (b) $\pi(2^n) < \dfrac{2^{n+1}\log_2(n-1)}{n}$; \item[(c)] Deduce that, for all real numbers $x \ge 8$, \[\pi(x) < \frac{4x\log_2(\log_2(x))}{\log_2(x)}.\] \end{enumerate} \end{enumerate}

\end{enumerate}

\end{document}
